\documentclass[a4paper,10pt]{article}
\usepackage[utf8]{inputenc}
\usepackage{graphicx}
\usepackage{subfigure}
\usepackage{fancyhdr}
\usepackage{listings}
\usepackage{ gensymb }


%opening

\pagestyle{fancy} %eigener Seitenstil
\fancyhf{} %alle Kopf- und Fußzeilenfelder bereinigen
\fancyhead[L]{SEP Abgabegespräch SS2015/16} %Kopfzeile links
\fancyhead[C]{} %zentrierte Kopfzeile
\fancyhead[R]{\includegraphics[width=1cm]{../../template/tugraz_bw}} %Kopfzeile rechts
\renewcommand{\headrulewidth}{0.4pt} %obere Trennlinie
\fancyfoot[C]{} %Seitennummer
\renewcommand{\footrulewidth}{0.4pt} %untere Trennlinie


\lstset{language=C++, basicstyle=\ttfamily\tiny}  

\begin{document}

\section*{Autos}

Schreibe ein Programm, das eine Liste von Autos aus einer Datei einliest. Jede Zeile beschreibt ein Auto und hat das Format \texttt{<brand> <speed> <color>}. Implementiere drei Funktionen:
\begin{description}
 \item[\texttt{void printCars(std::vector$<$Car$>\&$ cars)}] Bekommt eine Referenz auf einen Vektor von \texttt{Car} Objekten und gibt alle Autos in der Form \\
 \texttt{<brand>: <speed> (<color>)$\setminus$n} aus. 
 
   \textbf{Beispiel}  {\footnotesize\begin{verbatim}Audi: 280 (black)
Audi 280 silver
VW 220 silver\end{verbatim}}
 
 
 \item[\texttt{void sortCars(std::vector$<$Car$>\&$ cars)}] Bekommt eine Referenz auf einen Vektor von \texttt{Car} Objekten und sortiert diese absteigend nach Geschwindigkeit.
 
 \textbf{Beispiel} {\footnotesize\begin{verbatim}Audi: 330 (red)
BMW: 310 (white) 
BMW: 310 (black)\end{verbatim}}
 
 \item[\texttt{void removeDuplicateCars(std::vector$<$Car$>\&$ cars)}] Bekommt eine Referenz auf einen Vektor von \texttt{Car} Objekten und entfernt alle doppelten Autos. Autos sind gleich, wenn sowohl die Marke (\texttt{brand}) als auch die Geschwindigkeit (\texttt{speed}) übereinstimmen. (Es ist in Ordnung, wenn die Autos dabei auch sortiert werden.)
 
 \textbf{Beispiel} {\footnotesize\begin{verbatim}Audi: 330 (red)
BMW: 310 (white)
Audi: 280 (black)\end{verbatim}}
 
\end{description}

Ist die Anzahl der Parameter falsch, soll ``\texttt{Usage: ./cars <car file>}'' ausgegeben werden. Sollte eine Zeile in der Datei ungültig sein, soll sie ignoriert werden.

\vspace{1em}

Beispiel für eine Auto-Liste \texttt{cars.txt}:
\begin{lstlisting}[frame=single]
Audi 280 black
Audi 280 silver
BMW 310 white
VW 220 silver
Audi 330 red
VW 180 blue
BMW 310 black
VW 280 silver
\end{lstlisting}


\subsection*{Aufruf}
\texttt{./cars <car list>}

\subsection*{Bewertung}
\begin{itemize}
 \item Dokumentation und Programmierstil: 1 Punkt
 \item Einhalten der OOP Konzepte: 1 Punkt
 \item Struktur \& Korrektheit: 1 Punkt
\end{itemize}

\newpage


\end{document}
